% Real slides: «Sortieren» (Algorithmik, Kapitel 3), ported 1:1 from references/sorting_2025.odp.
% Step sequences of the original (one Impress slide per step) are one frame with overlays here.
% The card values are mapped to the cards in images/playing_cards/ (5→2, 6→5, 7→6, 8→7; 9, 10 stay).
% Input by real-d.tex or real-b.tex, which define \RTitle, \RSection and \RTask.
\slidesubject{Algorithmik}
\slidechapter{3}
\slidecoverimages{cover/algorithmik/playing-cards.jpg, cover/algorithmik/pencils.jpg, cover/algorithmik/maze.jpg,
  cover/algorithmik/signpost.jpg, cover/algorithmik/bridges.jpg, cover/algorithmik/subway-map.jpg}
\title{Sortieren}
\subtitle{Selectionsort · Bubblesort · Mergesort}
\author{Florian Köhl}

% one step of a card sequence: \cardstep<overlay>{cards}{parts}{arrows}{caption}
\newcommand<>{\cardstep}[4]{\only#5{%
  \cardArray{name=c, items={#1}, parts={#2}, size=0.11}#3%
  \par\vspace{15mm}#4}}
% a row of cells for the merge figures: \cellrow{y}{group size}{gap}{values}{name}
\newcommand{\cellrow}[5]{%
  \foreach \v [count=\i from 0] in {#4} {%
    \pgfmathsetmacro{\xx}{\i+#3*floor(\i/#2)}%
    \node[cell, minimum size=\designCellSize-1.2mm] (#5\i) at (\xx,#1) {\v};}}
% \ptr{node}: pointer from below onto a cell
\newcommand{\ptr}[1]{\draw[pointer] ([yshift=-8mm]#1.south) -- ([yshift=-1mm]#1.south);}

\begin{document}

\RTitle

\begin{frame}[text=large]{Sortieren}
  \begin{itemize}
    \item Was heisst Sortieren?
    \item Was sortieren Menschen?
    \item Was sortieren Computer?
  \end{itemize}
\end{frame}

\begin{frame}[text=large]{Sortieren}
  \begin{itemize}
    \item Was ist eine aufsteigend sortierte Liste?
  \end{itemize}
  \pause
  \begin{concept}[Aufsteigend sortiert]
    Eine Liste ist \term{aufsteigend sortiert}, sofern jedes Element in der Liste kleiner oder gleich gross ist
    wie alle nachfolgenden Elemente.
  \end{concept}
\end{frame}

\begin{frame}{Sortiert?}
  \vspace*{6mm}
  \only<1>{\cardArray{items={2h,5h,6h,7h,9h,10h}, size=0.12}}%
  \only<2>{\cardArray{items={2h,6h,5h,9h,7h,10h}, size=0.12}}%
  \only<3>{\cardArray{items={10h,9h,7h,6h,5h,2h}, size=0.12}}
\end{frame}

\begin{frame}[text=large]{Sortieralgorithmen}
  \begin{itemize}
    \item Bedingungen: Wir dürfen pro Zug …
    \begin{itemize}
      \item maximal zwei Karten aufgedeckt haben
      \item zwei Elemente miteinander vergleichen
      \item zwei Elemente miteinander vertauschen
    \end{itemize}
    \item Entwickelt paarweise ein Vorgehen, um eine Folge von (verdeckten) Karten zu sortieren!
  \end{itemize}
\end{frame}

\begin{frame}[doit={zu zweit}]{Let's go!}
  \vspace*{8mm}
  \cardArray{items={back,back,back,back,back,back}, size=0.12}
\end{frame}

\RSection{1}{Selectionsort}{Wiederholtes Suchen des Minimums}

\begin{frame}[text=small]{Selectionsort}
  \begin{enumerate}
    \item Teile die Elemente in einen sortierten (\term{S}) und einen unsortierten (\term{U}) Teil auf.
      S ist zu Beginn leer, U umfasst alle Elemente.
    \item Suche das kleinste Element in U.
    \item Tausche es mit dem ersten Element von U, sofern es nicht bereits dort steht.
    \item S wird um das eben gefundene Element vergrössert, U wird ein Element kürzer.
    \item Falls U mehr als 1 Element hat, gehe zu Schritt 2.
    \item Ende: Die Elemente sind sortiert.
  \end{enumerate}
\end{frame}

\begin{frame}{Selectionsort in Aktion}
  \vspace*{6mm}
  \cardstep<1>{5h,2h,9h,6h,10h,7h}{1/6/U}{}{}
  \cardstep<2>{5h,2h,9h,6h,10h,7h}{1/6/U}{\pointerbelow[8mm]{c2}{min}}{}
  \cardstep<3>{5h,2h,9h,6h,10h,7h}{1/6/U}{\pointerbelow[8mm]{c2}{min}\swaparrow{c1}{c2}}{}
  \cardstep<4>{2h,5h,9h,6h,10h,7h}{1/1/S,2/6/U}{}{U hat mehr als 1 Element → wiederholen}
  \cardstep<5>{2h,5h,9h,6h,10h,7h}{1/1/S,2/6/U}{\pointerbelow[8mm]{c2}{min}}{(keine Vertauschung notwendig)}
  \cardstep<6>{2h,5h,9h,6h,10h,7h}{1/2/S,3/6/U}{}{U hat mehr als 1 Element → wiederholen …}
  \cardstep<7>{2h,5h,9h,6h,10h,7h}{1/2/S,3/6/U}{\pointerbelow[8mm]{c4}{min}\swaparrow{c3}{c4}}{}
  \cardstep<8>{2h,5h,6h,9h,10h,7h}{1/3/S,4/6/U}{\pointerbelow[8mm]{c6}{min}\swaparrow{c4}{c6}}{}
  \cardstep<9>{2h,5h,6h,7h,10h,9h}{1/4/S,5/6/U}{\pointerbelow[8mm]{c6}{min}\swaparrow{c5}{c6}}{}
  \cardstep<10>{2h,5h,6h,7h,9h,10h}{1/5/S,6/6/U}{}{U hat nur 1 Element → fertig!}
\end{frame}

\begin{frame}{Fragen zu Selectionsort?}
\end{frame}

\begin{frame}[text=large]{Effizienz}
  \begin{itemize}
    \item Was ist ein effizienter Algorithmus?
  \end{itemize}
\end{frame}

\begin{frame}[text=large]{Effizienz: Laufzeit}
  \begin{itemize}
    \item Wie «schnell» ist Selectionsort?
    \item Wie messen wir die Effizienz eines solchen Algorithmus?
    \begin{itemize}
      \item Zeitmessung?
      \item Alternativen?
    \end{itemize}
  \end{itemize}
\end{frame}

\begin{frame}[text=large]{Effizienz: Laufzeit}
  \begin{itemize}
    \item Zeitmessung nicht immer geeignet!
    \begin{itemize}
      \item Unterschiedliche Hardware
      \item Unterschiedliche Implementierung
      \item Unterschiedliche Inputs (Eingaben)
    \end{itemize}
    \item<2-> Grobe Idee: Anzahl \term{Operationen} zählen
    \begin{itemize}
      \item<2-> «Wie viele Schritte braucht es?»
      \item<2-> Vergleiche zwischen zwei Elementen
      \item<2-> Vertauschungen zweier Elemente
    \end{itemize}
  \end{itemize}
\end{frame}

\begin{frame}[text=large]{Effizienz (Laufzeit)}
  \begin{itemize}
    \item \textbf{Problem 1:} Wir kennen die Grösse des Inputs nicht.
    \begin{itemize}
      \item z.\,B. wie lang die Liste ist, die wir sortieren sollen.
    \end{itemize}
    \item \textbf{Lösung:} Wir geben die Effizienz in Abhängigkeit von der Länge des Inputs an.
    \begin{itemize}
      \item Oft als $n$ bezeichnet.
    \end{itemize}
  \end{itemize}
\end{frame}

\begin{frame}[text=large]{Effizienz (Laufzeit)}
  \begin{itemize}
    \item \textbf{Problem 2:} Wir kennen den Input nicht.
    \begin{itemize}
      \item z.\,B. wie die Liste aussieht, die wir sortieren sollen.
    \end{itemize}
    \item \textbf{Lösung:} Wir geben die Effizienz für den besten oder schlechtesten Fall an.
    \begin{itemize}
      \item Oft auch für den «durchschnittlichen» Fall
    \end{itemize}
  \end{itemize}
\end{frame}

\begin{frame}{Beispiel Selectionsort}
  \begin{itemize}
    \item Liste mit 6 Elementen
    \only<1>{%
    \item Betrachte den ersten Durchgang des Algorithmus.
    \begin{itemize}
      \item Wie viele Vergleiche benötigen wir (bis wir das Minimum gefunden haben)?
      \item Wie viele Vertauschungen braucht es dann?
    \end{itemize}
    \item Wie sieht es in den nächsten Durchgängen aus?}%
    \only<2>{%
    \begin{itemize}
      \item 1.~Durchgang: 5 Vergleiche, max.~1 Vertauschung
      \item 2.~Durchgang: 4 Vergleiche, max.~1 Vertauschung
      \item 3.~Durchgang: 3 Vergleiche, max.~1 Vertauschung
      \item 4.~Durchgang: 2 Vergleiche, max.~1 Vertauschung
      \item 5.~Durchgang: 1 Vergleich, max.~1 Vertauschung
    \end{itemize}
    \item $1 + 2 + 3 + 4 + 5 = 15$ Vergleiche, mit max.~5 Vertauschungen.}
  \end{itemize}
\end{frame}

\begin{frame}{Beispiel Selectionsort (2)}
  \begin{itemize}
    \item Liste mit $n$ Elementen
    \only<1>{\item Wie viele «Durchgänge» benötigen wir?}%
    \only<2>{%
    \begin{itemize}
      \item 1.~Durchgang: $n-1$ Vergleiche, max.~1 Vertauschung
      \item 2.~Durchgang: $n-2$ Vergleiche, max.~1 Vertauschung
      \item 3.~Durchgang: $n-3$ Vergleiche, max.~1 Vertauschung
      \item …
      \item $(n-1)$.~Durchgang: 1 Vergleich, max.~1 Vertauschung
    \end{itemize}
    \item $1 + 2 + 3 + \dots + (n-3) + (n-2) + (n-1)$ Vergleiche}
  \end{itemize}
\end{frame}

\begin{frame}[text=large]{«Kleiner Gauss»}
  \vspace*{8mm}
  \[ 1 + 2 + 3 + \dots + (x-2) + (x-1) + x \;=\; \frac{x \cdot (x+1)}{2} \]
  \vspace{4mm}
  \[ \sum_{i=1}^{x} i \;=\; \frac{x \cdot (x+1)}{2} \]
\end{frame}

\begin{frame}{Beispiel Selectionsort (2)}
  \begin{itemize}
    \item $1 + 2 + 3 + \dots + (n-3) + (n-2) + (n-1)$ Vergleiche
  \end{itemize}
  \[ 1 + 2 + 3 + \dots + (n-2) + (n-1) \;=\; \frac{(n-1) \cdot n}{2} \;=\; \frac{n^2-n}{2} \]
  \begin{itemize}
    \item $n-1$ Durchgänge mit je max.~1 Vertauschung: max.~$n-1$ Vertauschungen
  \end{itemize}
\end{frame}

\begin{frame}[text=large]{Selectionsort: Effizienz}
  \begin{columns}[T,onlytextwidth]
    \begin{column}{0.48\textwidth}
      \begin{itemize}
        \item \textbf{Bester Fall:} (Wann?)
        \begin{itemize}
          \item 0 Vertauschungen
          \item $\frac{n^2-n}{2}$ Vergleiche
        \end{itemize}
      \end{itemize}
    \end{column}
    \begin{column}{0.48\textwidth}
      \begin{itemize}
        \item \textbf{Schlechtester Fall:} (Wann?)
        \begin{itemize}
          \item $n-1$ Vertauschungen
          \item $\frac{n^2-n}{2}$ Vergleiche
        \end{itemize}
      \end{itemize}
    \end{column}
  \end{columns}
\end{frame}

\RSection{2}{Bubblesort}{Wiederholtes paarweises Vertauschen}

\begin{frame}{Sortieren: Bubblesort}
  \begin{itemize}
    \item Gehe die Liste von links nach rechts durch.
    \begin{itemize}
      \item Vergleiche jeweils das aktuelle Element mit seinem rechten Nachbarn (falls vorhanden).
      \item Falls die beiden Elemente in der falschen Reihenfolge sind, vertausche sie.
    \end{itemize}
    \item Falls du in diesem Durchgang eine Vertauschung vorgenommen hast, wiederhole den Vorgang.
    \item Ansonsten ist die Liste sortiert.
  \end{itemize}
\end{frame}

\begin{frame}{Bubblesort in Aktion}
  \vspace*{6mm}
  \cardstep<1>{5h,2h,9h,6h,10h,7h}{}{\pointerbelow{c1}{}\pointerbelow{c2}{}\swaparrow{c1}{c2}}{}
  \cardstep<2>{2h,5h,9h,6h,10h,7h}{}{\pointerbelow{c2}{}\pointerbelow{c3}{}}{}
  \cardstep<3>{2h,5h,9h,6h,10h,7h}{}{\pointerbelow{c3}{}\pointerbelow{c4}{}\swaparrow{c3}{c4}}{}
  \cardstep<4>{2h,5h,6h,9h,10h,7h}{}{\pointerbelow{c4}{}\pointerbelow{c5}{}}{}
  \cardstep<5>{2h,5h,6h,9h,10h,7h}{}{\pointerbelow{c5}{}\pointerbelow{c6}{}\swaparrow{c5}{c6}}{}
  \cardstep<6>{2h,5h,6h,9h,7h,10h}{}{}{%
    \begin{itemize}
      \item 1.~Durchgang fertig!
      \item Haben wir vertauscht? Ja → also nochmals!
    \end{itemize}}
\end{frame}

\begin{frame}[doit={zu zweit}, text=large]{Ihr seid dran!}
  \vfill
  \RTask{38–41}{Sortieren von Hand}{Löst die Aufgaben 38–41 im Skript.}
  \vfill
\end{frame}

\begin{frame}[text=large]{Zwischenfazit}
  \begin{itemize}
    \item Selectionsort und Bubblesort haben beide eine \term{quadratische Laufzeit}.
    \begin{itemize}
      \item Die Laufzeitfunktion beinhaltet ein $n^2$.
      \item Verdoppelung der Inputgrösse führt zu Vervierfachung der Laufzeit.
    \end{itemize}
    \item Wie ineffizient ist das?
  \end{itemize}
\end{frame}

\begin{frame}{«Quadratische Laufzeit»: Beispiel}
  \vspace*{2mm}
  \centering\includegraphics[width=0.92\textwidth]{real/pictures/inbox.png}
\end{frame}

\RSection{3}{Mergesort}{Divide and Conquer}

\begin{frame}[text=large]{Divide and Conquer}
  \begin{itemize}
    \item Ansatz zum Lösen von Problemen
    \item \term{Divide}
    \begin{itemize}
      \item Aufteilen in Teilprobleme
      \item so lange, bis sie trivial lösbar sind
    \end{itemize}
    \item \term{Conquer}
    \begin{itemize}
      \item Kombination der Teillösungen zur Lösung des ganzen Problems
    \end{itemize}
  \end{itemize}
\end{frame}

\begin{frame}[text=large]{Beispiel: Mergesort}
  \begin{itemize}
    \item Effizientes Sortierverfahren
    \item Basiert auf \term{Divide and Conquer}
  \end{itemize}
  \vspace{8mm}
  \pythonArray{items={2,8,5,3,9,4,1,7}, drawindexes=false}
\end{frame}

\begin{frame}{Mergesort: Aufteilen}
  \only<1-2>{Wir teilen die Liste in zwei Teillisten auf.}%
  \only<3->{Gleicher Ansatz für die beiden resultierenden Teillisten.}
  \par\vspace{3mm}\centering
  \begin{tikzpicture}[x=\designCellSize, y=\designCellSize]
    \cellrow{0}{8}{0}{2,8,5,3,9,4,1,7}{a}
    \begin{scope}[visible on=<2->, xshift=-0.5*\designCellSize]\cellrow{-1.35}{4}{1}{2,8,5,3,9,4,1,7}{b}\end{scope}
    \begin{scope}[visible on=<3->, xshift=-1.2*\designCellSize]\cellrow{-2.7}{2}{0.8}{2,8,5,3,9,4,1,7}{c}\end{scope}
    \begin{scope}[visible on=<4->, xshift=-2.1*\designCellSize]\cellrow{-4.05}{1}{0.6}{2,8,5,3,9,4,1,7}{d}\end{scope}
  \end{tikzpicture}
\end{frame}

\begin{frame}[text=large]{Mergesort: Conquer}
  \begin{itemize}
    \item Resultat: viele kleine Listen der Länge 1
    \begin{itemize}
      \item Diese sind trivialerweise bereits sortiert!
    \end{itemize}
    \item Nun müssen wir diese Teillisten wieder zu einem Resultat zusammenfügen.
    \begin{itemize}
      \item Das «Merge» in Mergesort
      \item Wir machen das paarweise.
    \end{itemize}
  \end{itemize}
\end{frame}

\begin{frame}{Mergesort: Merge}
  «Merge» von zwei sortierten Listen:
  \begin{enumerate}
    \item Falls beide Listen nicht leer sind:
    \begin{enumerate}
      \item Vergleiche das kleinste noch nicht gemergte Element der linken mit demjenigen der rechten Liste.
      \item Wähle das kleinere der beiden aus und füge es in die neue Liste hinzu.
      \item Gehe zu Schritt 1.
    \end{enumerate}
    \item Falls eine der beiden Listen leer ist:
    \begin{enumerate}
      \item Füge alle Elemente der nicht leeren Liste der Reihe nach in die neue Liste ein.
    \end{enumerate}
  \end{enumerate}
\end{frame}

\begin{frame}{Mergesort: Conquer}
  \centering
  \begin{tikzpicture}[x=\designCellSize, y=\designCellSize]
    \cellrow{0}{1}{0.6}{2,8,5,3,9,4,1,7}{a}
    \begin{scope}[xshift=0.6*\designCellSize]\cellrow{-1.4}{2}{1.2}{2,8,3,5,4,9,1,7}{p}\end{scope}
    \foreach \i in {0,...,3} \node[cell blank, minimum size=\designCellSize-1.2mm] (m\i) at (\i+0.6,-3.6) {};
    \foreach \v [count=\i from 0] in {2,3,5,8} {\pgfmathtruncatemacro{\s}{\i+1}\node[cell done, visible on=<\s->, minimum size=\designCellSize-1.2mm] at (m\i) {\v};}
    \only<1>{\ptr{p0}\ptr{p2}}\only<2>{\ptr{p1}\ptr{p2}}\only<3>{\ptr{p1}\ptr{p3}}\only<4>{\ptr{p1}}
  \end{tikzpicture}

  \vspace{4mm}
  \only<1>{2 < 3: die 2 kommt in die Liste.}\only<2>{3 < 8: die 3 kommt in die Liste.}%
  \only<3>{5 < 8: die 5 kommt in die Liste.}\only<4>{Die rechte Liste ist leer: die 8 kommt in die Liste.}
  \vfill\stepdots{4}
\end{frame}

\begin{frame}[text=large]{Mergesort: Conquer}
  \begin{itemize}
    \item Gleiches Verfahren mit der rechten Teilliste.
    \item Dann bleiben noch 2 Teillisten (à je 4 Elemente) übrig.
  \end{itemize}
\end{frame}

\begin{frame}{Mergesort: Conquer}
  \centering
  \begin{tikzpicture}[x=\designCellSize, y=\designCellSize]
    \cellrow{0}{4}{1.4}{2,3,5,8,1,4,7,9}{p}
    \foreach \i in {0,...,7} \node[cell blank, minimum size=\designCellSize-1.2mm] (m\i) at (\i+0.7,-2.3) {};
    \foreach \v [count=\i from 0] in {1,2,3,4,5,7,8,9} {\pgfmathtruncatemacro{\s}{\i+1}\node[cell done, visible on=<\s->, minimum size=\designCellSize-1.2mm] at (m\i) {\v};}
    \only<1>{\ptr{p0}\ptr{p4}}\only<2>{\ptr{p0}\ptr{p5}}\only<3>{\ptr{p1}\ptr{p5}}\only<4>{\ptr{p2}\ptr{p5}}%
    \only<5>{\ptr{p2}\ptr{p6}}\only<6>{\ptr{p3}\ptr{p6}}\only<7>{\ptr{p3}\ptr{p7}}\only<8>{\ptr{p7}}
  \end{tikzpicture}

  \vspace{4mm}
  \only<1>{1 < 2: die 1 kommt in die Liste.}\only<2>{2 < 4: die 2 kommt in die Liste.}%
  \only<3>{3 < 4: die 3 kommt in die Liste.}\only<4>{4 < 5: die 4 kommt in die Liste.}%
  \only<5>{5 < 7: die 5 kommt in die Liste.}\only<6>{7 < 8: die 7 kommt in die Liste.}%
  \only<7>{8 < 9: die 8 kommt in die Liste.}\only<8>{Die linke Liste ist leer: die 9 kommt in die Liste.}
  \vfill\stepdots{8}
\end{frame}

\begin{frame}[text=large]{Mergesort: Effizienz}
  \begin{itemize}
    \item $O(n \cdot \log(n))$
    \item Sprich: Grössenordnung $n \cdot \log(n)$
    \item Ist das viel besser als «quadratisch»?
  \end{itemize}
\end{frame}

\begin{frame}{$n^2$ vs.\ $n \cdot \log(n)$}
  \centering
  \begin{tikzpicture}
    \begin{axis}[width=0.8\textwidth, height=0.52\textheight, axis lines=left, xmin=0, xmax=40, ymin=0, ymax=1600,
        xlabel={$n$}, axis line style={draw=mutedFg}, tick style={draw=mutedFg}, ticklabel style={font=\footnotesize, text=mutedFg},
        label style={text=textFg}, legend style={draw=none, fill=none, text=textFg, font=\footnotesize}, legend pos=north west,
        grid=major, grid style={draw=edgeCol!40!pageBg}]
      \addplot[termColor, line width=1.5pt, domain=1:40, samples=60] {x^2};
      \addplot[edgeCol, line width=1.5pt, domain=1:40, samples=60] {x*ln(x)/ln(2)};
      \legend{$n^2$, $n \cdot \log_2(n)$}
    \end{axis}
  \end{tikzpicture}
\end{frame}

\end{document}
